\documentclass[10pt,a4paper]{amsart}
\usepackage[margin=19mm]{geometry}
\usepackage{amsmath,amssymb,mathtools}
\usepackage[scr=rsfs,cal=euler]{mathalfa}
\usepackage{xfrac}
\usepackage[unicode,hidelinks,bookmarks=false]{hyperref}
\newcommand{\ie}{i.e.\ }
\newcommand{\cf}{cf.\ }
\newcommand{\resp}{respectively\ }
\newcommand{\Subsection}{\subsection}
\providecommand{\nobreakdash}{\nobreak\hbox{-}}
\setlength{\emergencystretch}{2em}
\multlinegap0pt
\usepackage{amssymb}
\usepackage{mathtools}
\usepackage{tcolorbox}
\usepackage{enumerate}
\usepackage{xcolor}
\usepackage{tikz}
\newtheorem{theorem}{Theorem}
\newcommand{\MF}{\mathrm{MF}}
\newcommand{\Veff}{V^{(\mathrm{eff})}}
\newcommand{\dd}{\mathrm{d}}
\newcommand{\pin}{p^{(\mathrm{in})}}
\newcommand{\ptar}{p^{(\mathrm{tar})}}
\newcommand{\tr}{\operatorname{Tr}}
\title{\textbf{Mean-Field Path-Integral Diffusion:\\ From Samples to Interacting Agents}}
\author{Michael Chertkov}
\date{Source: arXiv:2605.00007v1 (CC BY 4.0)}
\begin{document}
\maketitle

\noindent\emph{Source and licence.} \textbf{Mean-Field Path-Integral Diffusion:\\ From Samples to Interacting Agents}. Version arXiv:2605.00007v1 (CC BY 4.0), distributed by arXiv under the Creative Commons Attribution 4.0 International licence, http://creativecommons.org/licenses/by/4.0/, as recorded in the arXiv licence metadata for this exact version and shown on the abstract page https://arxiv.org/abs/2605.00007v1. Full licence notice: \texttt{licenses/arxiv-2605.00007-CC-BY-4.0.txt}.\par\smallskip

\noindent\textit{Editorial scope.} Counted unit: the article's theorem thm:linear-guidance (source0.tex lines 1299--1309). The article's complete original proof of it is delivered with it (source0.tex lines 1311--1369, one \texttt{proof} environment of 59 lines). No mathematics is written, repaired or re-derived: the statement and the proof are the article's own lines, in the article's own notation, and no retained line is substituted or re-flowed.  Retained uncounted as the source-local context the proof needs: lines 983--991; lines 1264--1267; lines 1269--1272; lines 1278--1281.  Nothing was omitted from any retained span, so the package carries no omission record.  No retained line contains a citation, so the wrapper carries no bibliography and no bibliography entry.  No retained line contains a figure environment, an image-inclusion command or drawing code, so no image is delivered and the package declares no asset dependency.  Editorial formatting only, in addition: an AMS \texttt{amsart} preamble that restates the article's own declarations and macros used by the retained lines, namely the environments \texttt{theorem} on separate counters, together with the standard packages those lines need.  The retained environments print in this document as Theorem 1 in order of appearance, whereas the article prints the same items with the numbering of its own full document.  The complete native source of the article as distributed in the arXiv e-print for this exact version is bundled unchanged as \texttt{sources/199531/source0.tex}.
\begin{equation}
  -\partial_t J_t
  = \underbrace{\int V_t(x-y)\,p_t(y)\,\dd y}_{\Veff_t(x)}
    + f_t(x)\cdot\nabla_x J_t
    + \tfrac{1}{2}\Delta_x J_t
    - \tfrac{1}{2}\|\nabla_x J_t\|^2,
  \quad J_1 = \varphi.
  \label{SI:eq:hjb-mf}
\end{equation}

\begin{equation}
  V_t(x) = \frac{\beta_t}{2}\|x\|^2, \quad \beta_t > 0,
  \label{SI:eq:V-quadratic}
\end{equation}

\begin{equation}
  \Veff_t(x) = \int V_t(x-y)\,p_t(y)\,\dd y= \frac{\beta_t}{2}\|x - m_t\|^2 + \frac{\beta_t}{2}\tr\Sigma_t,
  \label{SI:eq:Veff-quad}
\end{equation}

\begin{equation}
  \nu_t^{(\MF)}= \mathbb{E}_{x\sim p_t^{*}(\,\cdot\,|\,\nu^{(\MF)})}[x], \quad t\in[0,1].
  \label{SI:eq:self-consistency}
\end{equation}

\begin{theorem}[Linear MF guidance]
\label{thm:linear-guidance}
Let $f_t \doteq  0$, let $V_t$ be the quadratic potential \eqref{SI:eq:V-quadratic} with an arbitrary schedule $\beta_t > 0$, and let $\pin$, $\ptar$ be any probability measures with finite first moments $\bar m^{(\mathrm{in})} = \mathbb{E}_{\pin}[x]$ and $\bar m^{(\mathrm{tar})} = \mathbb{E}_{\ptar}[x]$. Also assume that the controlled process is initialized with $x_0\sim p^{(\mathrm{in})}$ (i.e. $p_0=p^{(\mathrm{in})}$), and that the MF fixed point exists and yields finite first moments  $m_t = \mathbb{E}_{p_t^*}[x]$ for all $t$. Then
\begin{equation}
  \nu_t^{(\MF)} = m_t
  = (1-t)\,\bar m^{(\mathrm{in})} + t\,\bar m^{(\mathrm{tar})}
  \quad \text{for all } t\in[0,1].
  \label{SI:eq:nu-linear}
\end{equation}
That is, the self-consistent MF guidance is the \emph{exact} linear interpolant between initial and target global means, independently of $\beta_t$ -- for any measurable $\beta_t>0$ such that the MF bridge is well-posed and $\mathbb E\|x_t\|<\infty$ for all $t$ -- and of the shape of $\pin$ and $\ptar$.
\end{theorem}

\begin{proof}
We show that $\ddot m_t = 0$.
\medskip

\noindent\textbf{Step 1 (It\^o + mean acceleration).}
With $f_t\doteq 0$, the McKean--Vlasov SDE is
$\dd x_t = u_t^*(x_t)\,\dd t + \dd W_t$.
Differentiating $m_t = \mathbb{E}[x_t]$ gives
$\dot m_t = \mathbb{E}[u_t^*(x_t)]$.
Applying It\^o's formula to $u_t^*(x_t)$ and taking expectations:
\begin{equation}
  \ddot m_t
  = \mathbb{E}\!\left[
      \partial_t u_t^*
      + (u_t^*\cdot\nabla_x)u_t^*
      + \tfrac{1}{2}\Delta u_t^*
    \right].
  \label{SI:eq:proof-ddot-m}
\end{equation}

\noindent\textbf{Step 2 (Differentiated HJB in space).} With $J_t = -\log\psi_t$ and $u_t^* = \nabla_x\log\psi_t$, the HJB equation \eqref{SI:eq:hjb-mf} with $f_t=0$ reads
\begin{equation*}
  \partial_t J_t= -\Veff_t + \tfrac{1}{2}|u_t^*|^2 + \tfrac{1}{2}\Delta\log\psi_t.
\end{equation*}
Taking the gradient in $x$ and using the symmetry of $\nabla_xu_t^* = \mathrm{Hess}(\log\psi_t)$:
\begin{equation}
  \partial_t u_t^*= \nabla_x\Veff_t- \underbrace{(\nabla_xu_t^*)\,u_t^*}_{=(u_t^*\cdot\nabla_x)u_t^*}- \tfrac{1}{2}\Delta u_t^*.
  \label{SI:eq:proof-dtu}
\end{equation}

\noindent\textbf{Step 3 (Cancellation).} Substituting \eqref{SI:eq:proof-dtu} into \eqref{SI:eq:proof-ddot-m}, the $(u_t^*\cdot\nabla_x)u_t^*$ and $\tfrac{1}{2}\Delta u_t^*$ terms cancel exactly, leaving
\begin{equation*}
  \ddot m_t = \mathbb{E}\!\left[\nabla_x\Veff_t(x_t)\right].
%  \label{SI:eq:proof-ddot-simple}
\end{equation*}
This identity holds for any interaction potential; the complex score structure of $u_t^*$ does not appear.

\noindent\textbf{Step 4 (Quadratic potential + self-consistency).} For the quadratic potential \eqref{SI:eq:Veff-quad}:
\begin{equation*}
  \nabla_x\Veff_t(x) = \beta_t(x - \nu_t^{(\MF)}).
\end{equation*}
Therefore
\begin{equation*}
  \ddot m_t
  = \beta_t\,\mathbb{E}[x_t - \nu_t^{(\MF)}]
  = \beta_t\,(m_t - \nu_t^{(\MF)}).
\end{equation*}
At the MF fixed point, the self-consistency condition
\eqref{SI:eq:self-consistency} gives $\nu_t^{(\MF)} = m_t$, so
\begin{equation*}
  \ddot m_t = \beta_t\,(m_t - m_t) = 0.
\end{equation*}

\noindent\textbf{Conclusion.}
With $\ddot m_t = 0$ and boundary conditions
$m_0 = \bar m^{(\mathrm{in})}$, $m_1 = \bar m^{(\mathrm{tar})}$,
we obtain $m_t = (1-t)\,\bar m^{(\mathrm{in})} + t\,\bar m^{(\mathrm{tar})}$,
and since $\nu_t^{(\MF)} = m_t$, the claim \eqref{SI:eq:nu-linear} follows.
\end{proof}

\end{document}
